\subsection*{1.3. 二次奇点的模}

以下命题由 R. Elkik 的定理 1.1.4 以及留给读者的显式计算推出(这些计算本质上已包含于 VI 6).

\subsubsection*{命题 1.3.1.}
\underline{设 $S=\operatorname{Spec}(A)$ 为闭点为 $s$ 的 Hensel 局部概形, $Y_0\subset\mathbb A_{k(s)}^{n+1}$ 为 1.2.3 或 1.2.4 所述类型的概形.}

\begin{enumerate}[label=(\roman*),leftmargin=2.6em]
\item \underline{假设 $Y_0$ 为 1.2.3 所述类型, 由方程}
\[
Q(\underline X)=\sum_{i\leq j}\bar a_{ij}x_i x_j=0
\qquad(\bar a_{ij}\in k(s)),
\]
\underline{定义, 并将 $\bar a_{ij}$ 提升为 $A$ 中的元素 $a_{ij}$. 则 $Y_0$ 在 $y_0$ 处奇点的形变的万有 Hensel $S$-概形为}
\[
T=\operatorname{Spec}(A\{b\}),
\]
\underline{其万有形变为由方程}
\[
\sum a_{ij}x_i x_j-b=0.
\]
\underline{定义的 $T$-概形 $Y\subset\mathbb A_T^{n+1}$ 在 $y_0$ 处的 Hensel 化.}

\item \underline{假设 $Y_0$ 为 1.2.4 所述类型, 由方程}
\[
Q(\underline x)=(x_0^2-\bar a)+\sum\bar a_{ij}x_i x_j=0
\qquad(\bar a,\bar a_{ij}\in k(s)),
\]
\underline{定义, 并将 $\bar a_{ij}$ 和 $\bar a$ 提升为 $A$ 中的元素 $a_{ij}$ 和 $a$. 则 $Y_0$ 在 $y_0$ 处奇点的形变的万有 Hensel $S$-概形为}
\[
T=\operatorname{Spec}(A\{b,c\}),
\]
\underline{其万有形变为由方程}
\[
(x_0^2-a)+\sum a_{ij}x_i x_j+b x_0+c=0.
\]
\underline{定义的 $T$-概形 $Y\subset\mathbb A_T^{n+1}$ 在 $y_0$ 处的 Hensel 化.}
\end{enumerate}

以下推论立即由 1.2.6 及 1.3.1 得出.

\subsubsection*{推论 1.3.2.}
\underline{设 $S=\operatorname{Spec}(A)$ 为闭点为 $s$ 的 Hensel 局部概形, $f:X\to S$ 为有限表现平坦态射. 假设在 $X_s$ 的闭点 $x$ 处, 概形 $X_s$ 有普通二次奇点, $k(x)$ 在 $k(s)$ 上纯不可分, 且 $X_s$ 的维数为 $n$.}

\begin{enumerate}[label=(\roman*),leftmargin=2.6em]
\item 若 $x$ 非退化, \underline{则存在系数属于 $A$ 的 $n+1$ 元非退化二次型 $Q$, 以及 $A$ 的极大理想中的元素 $b$, 使 $X$ 在 $x$ 处的 Hensel 化 $S$-概形同构于仿射空间 $\mathbb A_S^{n+1}$ 中由方程 $Q-b=0$ 定义的闭子概形在原点处的 Hensel 化, 如 1.3.1(i).}

\item 若 $x$ 退化($n=2m$, $\operatorname{char}(k(s))=2$), \underline{则存在系数属于 $A$ 的 $n+1$ 元非齐次二次型}
\begin{equation}\tag{1.3.2.1}
Q(\underline X)=x_0^2+b x_0+c+
 \sum_{1\leq i\leq j\leq 2m}a_{ij}x_i x_j,
\end{equation}
\underline{使 $b$ 属于极大理想, $2m$ 元二次型 $\sum a_{ij}x_i x_j$ 非退化, 且 $X$ 在 $x$ 处的 Hensel 化 $S$-概形同构于仿射空间 $\mathbb A_S^{n+1}$ 中由方程 $Q=0$ 定义的闭子概形在 $(\sqrt c,0,\ldots,0)$ 处的 Hensel 化.}
\footnote{印本将此公式标为 (4.14.2), 且首项写为 $x_0$. 上下文中的正规形及注 1.3.3 要求使用本节编号 (1.3.2.1) 和二次项 $x_0^2$; 英译均予恢复.}
\end{enumerate}

\subsubsection*{注 1.3.3.}
在情形 (i) 中, 假设 $X_s$ 在 $x$ 处的 Hensel 化经同构 $\varphi$ 同构于由方程
\[
\sum\bar a_{ij}x_i x_j=0,
\]
定义的闭子概形 $X_0\subset\mathbb A_k^{n+1}$ 在 $O$ 处的 Hensel 化, 并假设 $a_{ij}$ 是 $\bar a_{ij}$ 的提升. 此时可取
\[
Q=\sum a_{ij}x_i x_j.
\]

在情形 (ii) 中, 假设 $X_s$ 在 $x$ 处的 Hensel 化经同构 $\varphi$ 同构于由方程
\[
(x_0^2+\bar a)+\sum\bar a_{ij}x_i x_j=0,
\]
定义的闭子概形 $X_0\subset\mathbb A_k^{n+1}$ 在 $(\sqrt a,0,\ldots,0)$ 处的 Hensel 化, 并假设 $a_{ij}$ 是 $\bar a_{ij}$ 的提升. 此时可取
\[
Q=x_0^2+b x_0+c+\sum a_{ij}x_i x_j,
\qquad c\equiv\bar a\pmod{\mathfrak m_A}.
\]
在两种情形中, 都可选取 1.3.2 中保证存在的同构, 使它延拓 $\varphi$.

\subsubsection*{推论 1.3.4.}
\underline{设 $f:X\to S$ 如 1.3.2. 则 $x$ 在 $X$ 中存在一个邻域 $U$, 使包含于 $U$ 的 $f$ 的非光滑点都是其各自纤维的普通二次点.}